\subsection{Aplicación a los morfismos propios: II. Criterios de planitud cohomológica}
\label{subsection:III.7.8}
\label{III.7.8}

\begin{definition}[7.8.1]
\label{III.7.8.1}
Sean $X$ e $Y$ dos preesquemas, sea $f:X\to Y$ un morfismo
cuasicompacto y separado, sea $\mathscr P_\bullet$ un complejo de
$\mathscr O_X$-módulos cuasicoherentes y planos sobre $Y$, y sea
$\mathscr T_\bullet$ el funtor homológico de $\mathscr O_Y$-módulos
cuasicoherentes definido por \eqref{III.7.7.1.1}. Sea $y$ un punto
de $Y$.

Se dice que $\mathscr P_\bullet$ es
\oldpage[III]{205}
\emph{homológicamente plano sobre
$Y$ en el punto $y$, en dimensión $p$} (o \emph{cohomológicamente
plano sobre $Y$ en el punto $y$, en dimensión $-p$}) si existe un
entorno abierto $U$ de $y$ en $Y$ tal que
\[
  \mathscr T_p^U=\mathscr T_U^{-p}
\]
es exacto.

Se dice que $\mathscr P_\bullet$ es \emph{homológicamente plano en
dimensión $p$ sobre $Y$} (o \emph{cohomológicamente plano en
dimensión $-p$ sobre $Y$}) si es homológicamente plano sobre $Y$ en
todo punto $y\in Y$, en dimensión $p$.

Cuando $\mathscr P_\bullet$ es homológicamente plano sobre $Y$
(respectivamente, sobre $Y$ en el punto $y$) en \emph{toda}
dimensión $p$, se dice simplemente que $\mathscr P_\bullet$ es
\emph{homológicamente plano sobre $Y$} (respectivamente, sobre $Y$
en el punto $y$), o \emph{cohomológicamente plano sobre $Y$}
(respectivamente, sobre $Y$ en el punto $y$).
\end{definition}

\begin{env}[7.8.2]
\label{III.7.8.2}
Por definición, la noción de planitud homológica sobre $Y$ es
\emph{local sobre $Y$}.

Si $Y$ es localmente noetheriano, o si $Y$ es un esquema, para que
$\mathscr P_\bullet$ sea homológicamente plano sobre $Y$ en
dimensión $p$, es necesario y suficiente que el funtor
$\mathscr T_p$ sea exacto. La demostración para $Y$ localmente
noetheriano se dio durante la demostración de \sref{III.7.7.5};
cuando $Y$ es un esquema, el razonamiento es el mismo, basado en
\sref[I]{I.9.4.2} aplicado a un abierto afín de un esquema
cuasicompacto.
\end{env}

\begin{proposition}[7.8.3]
\label{III.7.8.3}
Con las notaciones e hipótesis de \sref{III.7.8.1}, las condiciones
siguientes son equivalentes:
\begin{enumerate}
  \item[a)] $\mathscr P_\bullet$ es homológicamente plano sobre $Y$
    en el punto $y$, en dimensión $p$.
  \item[b)] Existe un entorno abierto $U$ de $y$ en $Y$ tal que
    $\mathscr T_p^U$ y $\mathscr T_{p+1}^U$ son exactos por la
    derecha.
  \item[c)] Existe un entorno abierto $U$ de $y$ en $Y$ tal que
    $\mathscr T_p^U$ y $\mathscr T_{p-1}^U$ son exactos por la
    izquierda.
  \item[d)] Existe un entorno abierto $U$ de $y$ en $Y$ tal que
    $\mathscr T_{p+1}^U$ es exacto por la derecha y
    $\mathscr T_{p-1}^U$ es exacto por la izquierda.
\end{enumerate}
\end{proposition}

\begin{proof}
Teniendo en cuenta la interpretación de $\mathscr T_p$ cuando $Y$
es afín, esto no es más que una traducción de una parte de
\sref{III.7.3.3}.
\end{proof}

\begin{proposition}[7.8.4]
\label{III.7.8.4}
Sean $Y$ un preesquema localmente noetheriano, $f:X\to Y$ un
morfismo propio, $\mathscr P_\bullet$ un complejo acotado
inferiormente de $\mathscr O_X$-módulos coherentes y planos sobre
$Y$, y $\mathscr T_\bullet$ el funtor definido por
\eqref{III.7.7.1.1}. Para todo $y\in Y$, las condiciones siguientes
son equivalentes:
\begin{enumerate}
  \item[a)] $\mathscr P_\bullet$ es homológicamente plano sobre $Y$
    en $y$, en dimensión $p$.
  \item[b)] El funtor $\mathscr T_p^{\mathscr O_y}$ es exacto.
  \item[c)] Existe un entero $n_0$ tal que, para $n\geq n_0$,
    \[
      \operatorname{length}
        \mathscr T_p^{\mathscr O_y}
          \bigl(\mathscr O_y/\mathfrak m_y^{\,n+1}\bigr)
      =
      \operatorname{length}
        \mathscr T_p^{\mathscr O_y}\bigl(\kres(y)\bigr)
      \cdot
      \operatorname{length}
        \bigl(\mathscr O_y/\mathfrak m_y^{\,n+1}\bigr).
      \tag{7.8.4.1}\label{III.7.8.4.1}
    \]
    Aquí todas las longitudes son longitudes de
    $\mathscr O_y$-módulos.
  \item[d)] Existe un entorno abierto $U$ de $y$ tal que
    \[
      \mathscr H^{-p}(f,\mathscr P^\bullet)|U
    \]
    es isomorfo a un $(\mathscr O_Y|U)$-módulo de la forma
    $(\mathscr O_Y|U)^m$ y, para todo $(\mathscr O_Y|U)$-módulo
    cuasicoherente $\mathscr M$, el homomorfismo canónico
    \[
      \mathscr H^{-p}(f,\mathscr P^\bullet)|U
        \otimes_{\mathscr O_Y|U}\mathscr M
      \longrightarrow
      \mathscr H^{-p}
        \bigl(f,(\mathscr P^\bullet|U)
          \otimes_{\mathscr O_Y|U}\mathscr M\bigr)
      \tag{7.8.4.2}\label{III.7.8.4.2}
    \]
    es biyectivo.
\end{enumerate}
Cuando se cumplen estas condiciones, se tiene además:
\begin{enumerate}
  \item[e)] Existe un entorno de $y$ en el cual la función
    $z\mapsto d_p(z)$, definida en \eqref{III.7.7.5.1}, es
    constante.
\end{enumerate}
Además, si $Y$ es reducido en $y$ \sref[0]{0.4.1.4}, la condición
e) es equivalente a las demás condiciones.
\end{proposition}

\oldpage[III]{206}
\begin{proof}
En efecto, la condición b) equivale a decir que
$\mathscr T_p^{\mathscr O_y}$ y $\mathscr T_{p+1}^{\mathscr O_y}$
son exactos por la derecha \sref{III.7.3.3}. La equivalencia de a)
y b) resulta entonces de \sref{III.7.7.10} y \sref{III.7.8.3}.

Como $\mathscr O_y/\mathfrak m_y^{\,n+1}$ es artiniano, y
\[
  \mathscr T_p^{\mathscr O_y}
    \bigl(\mathscr O_y/\mathfrak m_y^{\,n+1}\bigr)
  \quad\hbox{y}\quad
  \mathscr T_p^{\mathscr O_y}\bigl(\kres(y)\bigr)
\]
son $(\mathscr O_y/\mathfrak m_y^{\,n+1})$-módulos de tipo finito
\sref{III.7.7.4}, y por tanto de longitud finita, la equivalencia de
b) y c) resulta nuevamente de \sref{III.7.7.10} y
\sref{III.7.3.5.7}.

Que a) implica e), y que a) y e) son equivalentes cuando
$\mathscr O_y$ es reducido, resulta de \sref{III.7.6.9}.
Finalmente, a) implica que \eqref{III.7.8.4.2} es biyectivo, por la
definición \sref{III.7.8.1} y \sref{III.7.7.5}. También implica que
\[
  \bigl(\mathscr H^{-p}(f,\mathscr P^\bullet)\bigr)_y
\]
es un $\mathscr O_y$-módulo plano por la condición (f) de
\sref{III.7.3.3}, y por tanto un $\mathscr O_y$-módulo libre de
tipo finito \sref[0]{0.10.1.3}, pues es de tipo finito por
\sref{III.7.7.4}. Como $\mathscr H^{-p}(f,\mathscr P^\bullet)$ es
un $\mathscr O_Y$-módulo coherente \sref{III.7.7.4}, es localmente
libre en un entorno de $y$ \sref[0]{0.5.2.7}. Recíprocamente, d)
implica claramente a), por la definición del funtor
$\mathscr T_p^U$ \eqref{III.7.7.2.2}.
\end{proof}

\begin{proposition}[7.8.5]
\label{III.7.8.5}
Bajo las hipótesis de \sref{III.7.8.4}, las condiciones siguientes
son equivalentes:
\begin{enumerate}
  \item[a)] $\mathscr P_\bullet$ es homológicamente plano sobre $Y$
    en toda dimensión $i\leq p$.
  \item[b)] Para $i\leq p+1$, los funtores $\mathscr T_i$ son
    exactos por la derecha.
  \item[c)] Para $i\geq-p$, los $\mathscr O_Y$-módulos
    $\mathscr H^i(f,\mathscr P^\bullet)$ son localmente libres.
\end{enumerate}
\end{proposition}

\begin{proof}
La equivalencia de a) y b) es inmediata por \sref{III.7.8.3}, y a)
implica c) por \sref{III.7.8.4}.

Recíprocamente, supongamos que se cumple c). Recordemos también que
$\mathscr L_i=0$ para $i<i_0$ \sref{III.7.7.4}, y por tanto
$\mathscr T_i=0$ para $i\leq i_0$. Todo punto $y\in Y$ tiene, pues,
un entorno afín $U=\Spec(A)$ tal que $T_i^A(A)$ es un $A$-módulo
libre para $i\leq p$. De \sref{III.7.3.7} se deduce que
$\mathscr T_i^U=T_i^A$ es exacto para $i\leq p$.
\end{proof}

Aplicaremos principalmente los criterios de planitud cohomológica al
caso en que el complejo $\mathscr P_\bullet$ está concentrado en un
\emph{solo} $\mathscr O_X$-módulo coherente $\mathscr F$, plano
sobre $Y$, tomado como $\mathscr P_0$. Recordemos que en este caso
\[
  \mathscr T_p(\mathscr M)
  =R^{-p}f_*
    \bigl(\mathscr F\otimes_{\mathscr O_Y}\mathscr M\bigr).
\]

\begin{proposition}[7.8.6]
\label{III.7.8.6}
Sean $Y$ un preesquema localmente noetheriano, $f:X\to Y$ un
morfismo propio y plano, y $y$ un punto de $Y$. Denotemos por $X_y$
la fibra $f^{-1}(y)=X\otimes_Y\kres(y)$. Supongamos que
\[
  \Gamma(X_y,\mathscr O_{X_y})=R
\]
es una $\kres(y)$-álgebra separable (Bourbaki,
\emph{Algèbre}, cap.~VIII, \S~7, n\textsuperscript{o}~5), es decir,
un producto directo de un número finito de extensiones separables de
grado finito de $\kres(y)$. Entonces $\mathscr O_X$ es
cohomológicamente plano sobre $Y$ en el punto $y$, en dimensión $0$.
\end{proposition}

\begin{proof}
Por \sref{III.7.8.4}, podemos restringirnos al caso en que $Y$ es
el espectro del anillo local $A=\mathscr O_y$. La planitud de $f$
implica $\mathscr T_{-1}=0$, de modo que $T_0^A$ ya es exacto por la
izquierda y solo queda probar que es exacto por la derecha. Por la
condición c) de \sref{III.7.7.10}, podemos reducirnos incluso al
caso en que $A=\mathscr O_y$ es artiniano.

Sea $k'$ una extensión finita de $\kres(y)$ que descomponga $R$, de
modo que
\[
  R\otimes_{\kres(y)}k'
\]
es un producto directo de un número finito de cuerpos isomorfos a
$k'$. Existe un homomorfismo local de $A$ en un anillo local $A'$
tal que $A'$ es una $A$-álgebra libre finita y su cuerpo residual es
isomorfo a $k'$ \sref[0]{0.10.3.2}. Por \sref{III.7.6.1}, basta,
pues, probar que $T_0^{A'}$ es exacto por la derecha; en otras
palabras, podemos suponer que $R$ es un producto directo de $m$
cuerpos isomorfos a $\kres(y)$.

Utilizaremos el siguiente lema elemental.

\begin{lemma}[7.8.6.1]
\label{III.7.8.6.1}
Sea $Z$ un espacio localmente anillado. Para que $Z$ sea conexo,
\oldpage[III]{207}
es necesario y suficiente que el anillo $\Gamma(Z,\mathscr O_Z)$
no sea producto de dos anillos no nulos.
\end{lemma}

\begin{proof}
Es claro que, si $Z$ es la unión de dos abiertos no vacíos y
disjuntos $Z_1$ y $Z_2$, entonces $\Gamma(Z,\mathscr O_Z)$ es
isomorfo al producto de los dos anillos no nulos
$\Gamma(Z_1,\mathscr O_Z)$ y $\Gamma(Z_2,\mathscr O_Z)$.

Recíprocamente, decir que $\Gamma(Z,\mathscr O_Z)$ es tal producto
equivale a decir que contiene un idempotente $s$ distinto de $0$ y
de $1$. Para todo $z\in Z$, el germen $s_z$ es entonces un
idempotente del anillo local $\mathscr O_z$, y por tanto es igual a
$0$ o a $1$. El conjunto de los $z$ tales que $s_z=0$ es abierto.
Por otra parte, $s_z=1$ implica por definición que $s(z)\ne0$, y el
conjunto de los $z$ tales que $s_z=1$ es también abierto
\sref[0]{0.5.5.2}. De aquí se sigue la conclusión.
\end{proof}

Del lema se deduce que $X_y$ tiene exactamente $m$ componentes
conexas $X_i'$ y que
\[
  \Gamma(X_i',\mathscr O_{X_i'})=\kres(y)
\]
para todo $i$. Como $A$ es local y artiniano, su espectro consta de
un solo punto y, por consiguiente, $X$ y $X_y$ tienen el mismo
espacio subyacente. Así, $X$ tiene $m$ componentes conexas $X_i$
tales que
\[
  X_i'=X_i\otimes_Y\kres(y).
\]
Quedamos finalmente reducidos al caso $R=\kres(y)$. Por la
condición b) de \sref{III.7.7.10}, queda probar que el homomorfismo
canónico
\[
  \Gamma(X,\mathscr O_X)\longrightarrow
  \Gamma(X_y,\mathscr O_{X_y})
\]
es sobreyectivo. Esto es inmediato, pues el compuesto
\[
  \Gamma(Y,\mathscr O_Y)=A
  \longrightarrow\Gamma(X,\mathscr O_X)
  \longrightarrow\Gamma(X_y,\mathscr O_{X_y})=\kres(y)
\]
ya es sobreyectivo.
\end{proof}

\begin{corollary}[7.8.7]
\label{III.7.8.7}
Bajo las hipótesis de \sref{III.7.8.6}, existe un entorno abierto
$U$ de $y$ tal que:
\begin{enumerate}
  \item[(i)] $f_*(\mathscr O_X)|U$ es isomorfo a un
    $(\mathscr O_Y|U)$-módulo de la forma $(\mathscr O_Y|U)^m$;
  \item[(ii)] para todo $z\in U$, el homomorfismo canónico
    \[
      \bigl(f_*(\mathscr O_X)\bigr)_z
        \otimes_{\mathscr O_z}\kres(z)
      \longrightarrow
      \Gamma(X_z,\mathscr O_{X_z})
    \]
    es biyectivo.
\end{enumerate}
\end{corollary}

\begin{proof}
La aserción (i) resulta de \sref{III.7.8.6} y \sref{III.7.8.4}.
La aserción (ii) resulta de la exactitud de $\mathscr T_0^U$, para
una elección adecuada de $U$, y de \eqref{III.7.7.5.3}.
\end{proof}

\begin{corollary}[7.8.8]
\label{III.7.8.8}
Supongamos que se cumplen las hipótesis de \sref{III.7.8.6} y,
además, que
\[
  \Gamma(X_y,\mathscr O_{X_y})=\kres(y).
\]
Entonces existe un entorno abierto $U$ de $y$ tal que el
homomorfismo canónico
\[
  \mathscr O_Y|U\longrightarrow f_*(\mathscr O_X)|U
\]
es biyectivo.
\end{corollary}

\begin{proof}
En efecto, de la aserción (ii) de \sref{III.7.8.7} se deduce que el
entero $m$ que figura en la aserción (i) de ese corolario es
necesariamente igual a $1$.
\end{proof}

\begin{corollary}[7.8.9]
\label{III.7.8.9}
Bajo las hipótesis de \sref{III.7.8.6}, existen un entorno abierto
$U$ de $y$, un $\mathscr O_U$-módulo coherente $\mathscr Q$,
determinado salvo isomorfismo único, y un isomorfismo de funtores en
el $\mathscr O_U$-módulo cuasicoherente $\mathscr M$,
\[
  R^1f_*\bigl(f^*(\mathscr M)\bigr)
  \xrightarrow{\sim}
  \mathscr Hom_{\mathscr O_U}(\mathscr Q,\mathscr M).
  \tag{7.8.9.1}\label{III.7.8.9.1}
\]
\end{corollary}

\begin{proof}
La hipótesis implica que $\mathscr T_0^U$ es exacto para un $U$
adecuado. Basta, por tanto, aplicar la equivalencia de las
condiciones a) y b') de \sref{III.7.7.5}, con $p=0$ y tomando como
$\mathscr P_\bullet$ el complejo concentrado en grado $0$ cuyo
término es $\mathscr O_X$.
\end{proof}

\begin{remarks}[7.8.10]
\label{III.7.8.10}
\textnormal{(i)}
Bajo las hipótesis de \sref{III.7.8.6}, consideremos la
factorización de Stein de $f$ \sref{III.4.3.3},
\[
  \xymatrix{
    X \ar[r]^{f'} & Y' \ar[r]^{g} & Y .
  }
\]
\oldpage[III]{208}
Aquí
\[
  Y'=\Spec\bigl(f_*(\mathscr O_X)\bigr).
\]
El morfismo finito $g$ es tal que
\[
  g_*(\mathscr O_{Y'})=f_*(\mathscr O_X)
\]
es localmente libre cerca de $y$, y su fibra sobre $y$ es el
espectro de una álgebra separable sobre $\kres(y)$
\sref[II]{II.1.5.1}. Deduciremos en el capítulo~IV que existe un
entorno abierto $U$ de $y$ en $Y$ tal que, para la restricción
$g^{-1}(U)\to U$ de $g$, \emph{toda} fibra $g^{-1}(z)$, con
$z\in U$, es el espectro de una álgebra separable sobre $\kres(z)$.
Esto es lo que llamaremos un \emph{revestimiento étale} de $U$. Se
deducirá entonces de la aserción (ii) de \sref{III.7.8.7} que la
hipótesis sobre el punto $y$ en \sref{III.7.8.6} se cumple también
en todo punto de un entorno de $y$.

\medskip
\textnormal{(ii)}
Veremos en el capítulo~V que, incluso si $X$ es proyectivo sobre $Y$
(e incluso si además es ``simple'' sobre $Y$, propiedad que se
definirá en el capítulo~IV), el $\mathscr O_U$-módulo $\mathscr Q$
de \sref{III.7.8.9} no es necesariamente localmente libre. En otras
palabras, bajo estas hipótesis $\mathscr O_X$ no es necesariamente
cohomológicamente plano en \emph{dimensión $1$} sobre $Y$ en el
punto $y$. En el capítulo~V interpretaremos $\mathscr Q$ como el haz
de $1$-formas diferenciales del esquema de Picard de $X$ relativo a
$Y$, a lo largo de la sección unidad.
\end{remarks}
